%% =====================================================================
%%  Corrigés — Unit 5 « What's this? »
%% =====================================================================
\corrigesUnite{5}{What's this?}

\begin{corrige}{1}
\begin{reponses}[4]
\item \rep{pens}
\item \rep{boxes}
\item \rep{watches}
\item \rep{cities}
\item \rep{keys}
\item \rep{children}
\item \rep{men}
\item \rep{women}
\end{reponses}
\pourquoi{\en{city} : consonne + \en{y} → \en{cities} ; mais \en{key} : voyelle + \en{y} → \en{keys}.
\en{child}, \en{man} et \en{woman} ont des pluriels irréguliers, à apprendre par cœur.}
\end{corrige}

\begin{corrige}{2}
\begin{reponses}[4]
\item \rep{a}
\item \rep{The}
\item \rep{an}
\item \rep{the}
\item \rep{an}
\item \rep{The}
\item \rep{the}
\item \rep{a}
\end{reponses}
\pourquoi{on dit \en{a} / \en{an} la première fois qu'on parle d'une chose, puis \en{the} quand elle
est connue : \en{It's a cat. The cat is black.} Au numéro 4, il n'y a qu'une machine à café :
\en{the}. Londres est « une grande ville » parmi d'autres : \en{a}.}
\end{corrige}

\begin{corrige}{3}
\begin{reponses}[3]
\item \rep{This}
\item \rep{Those}
\item \rep{These}
\item \rep{That}
\item \rep{this}
\item \rep{those}
\end{reponses}
\retiens{près : \en{this} (singulier), \en{these} (pluriel) ; loin : \en{that}, \en{those}.}
\end{corrige}

\begin{corrige}{4}
\begin{reponses}[3]
\item \rep{They're keys.}
\item \rep{It's an umbrella.}
\item \rep{They're books.}
\item \rep{It's a lamp.}
\item \rep{They're watches.}
\item \rep{It's a cup.}
\end{reponses}
\piege{\en{It's a umbrella} : devant le son voyelle de \en{umbrella}, il faut \en{an}.}
\end{corrige}

\begin{corrige}{5}
\begin{reponses}[3]
\item \rep{blue}
\item \rep{white}
\item \rep{yellow}
\item \rep{red}
\item \rep{brown}
\item \rep{green}
\end{reponses}
\end{corrige}

\begin{corrige}{6}
\begin{reponses}[1]
\item \rep{I have two blue pens.}
\item \rep{It's an umbrella.}
\item \rep{These books are new.}
\item \rep{These keys are red.}
\item \rep{That's my bag.}
\item \rep{The three children are in London.}
\end{reponses}
\pourquoi{l'adjectif ne prend jamais de \en{-s} (phrase 1) ; avec un nom pluriel, \en{this} devient
\en{these} (phrase 4) ; \en{That's} = \en{That is}, au singulier (phrase 5).}
\end{corrige}

\begin{corrige}{7}
\begin{reponses}[3]
\item \rep{/s/}
\item \rep{/z/}
\item \rep{/ɪz/}
\item \rep{/s/}
\item \rep{/z/}
\item \rep{/ɪz/}
\item \rep{/s/}
\item \rep{/z/}
\item \rep{/ɪz/}
\end{reponses}
\pourquoi{\phon{s} après \phon{p} et \phon{k} (\en{cups}, \en{desks}, \en{books}) ; \phon{z} après un
son voisé (\en{bags}, \en{phones}, \en{keys}) ; \phon{ɪz} quand on ajoute une syllabe
(\en{boxes}, \en{watches}, \en{glasses}). Texte de la piste page~\pageref{trans-5.10}.}
\end{corrige}

\begin{corrige}{8}
\begin{reponses}[2]
\item \rep{Emma}
\item \rep{pink}
\item \rep{Léa}
\item \rep{green}
\item \rep{Mr. Carter}
\item \rep{black}
\item \rep{Maya}
\item \rep{purple}
\end{reponses}
\pourquoi{repérez les génitifs et les possessifs : \en{It's Emma's}, \en{his umbrella} (celui de
Mr.\ Carter), \en{It's Maya's}.}
\end{corrige}

\begin{corrige}{9}
\begin{anglais}\otTableau
\begin{tabular}{@{}l@{\hspace{6mm}}l@{}}
\otNomDit{sam} & What's \rep{this}?\\
\otNomDit{lea} & It's \rep{a} lamp from Lyon.\\
\otNomDit{sam} & And what are \rep{these}?\\
\otNomDit{lea} & They're my \rep{books}.\\
\otNomDit{sam} & Are \rep{those} your plants?\\
\otNomDit{lea} & Yes, \rep{they} are. They're from Mrs.\ Lopez.\\
\otNomDit{sam} & \rep{The} lamp is orange!\\
\otNomDit{lea} & Yes, and all my \rep{boxes} are brown!\\
\end{tabular}
\end{anglais}
\retiens{\en{all my boxes} = « tous mes cartons ». \en{all} se place devant le possessif.}
\end{corrige}

\begin{corrige}{10}
\begin{reponses}[1]
\item \rep{What's this? — It's an umbrella.}
\item \rep{They're my keys.} \textit{ou} \rep{These are my keys.}
\item \rep{This book is blue.}
\item \rep{Those pens are Jack's.}
\item \rep{The children are in London.}
\item \rep{It's an old black bag.}
\item \rep{What color is your lamp? — It's orange.}
\item \rep{These boxes are small.}
\end{reponses}
\pourquoi{phrase 6 : l'ordre est « article + adjectifs + nom », et \en{an} car \en{old} commence par
une voyelle. Phrase 4 : \en{Jack's} suffit pour dire « à Jack ».}
\end{corrige}

\begin{corrige}{11}
Réponse modèle :
\begin{anglais}
\begin{tcolorbox}[otBase,colback=otGreenTint,colframe=otGreenTint,fontupper=\fontsize{10.4}{14}\selectfont]
This is my laptop. It's gray and it's new. These are my pens: they're blue and black. This is my
phone. It's pink! That's my lamp. It's white. And those are my books. They're old.
\end{tcolorbox}
\end{anglais}
\retiens{\en{This is…} / \en{That's…} pour un objet, \en{These are…} / \en{Those are…} pour plusieurs ;
la couleur reste au singulier.}
\end{corrige}

\begin{corrige}{12}
Réponses modèles :
\begin{reponses}[1]
\item \rep{No, it isn't. It's Jack's.}
\item \rep{It's black.} \textit{ou} \rep{My bag is brown.}
\item \rep{They're my keys.}
\item \rep{Yes, it is.}
\item \rep{No, it isn't. It's old.}
\end{reponses}
\end{corrige}

\begin{corrigetest}
\begin{reponses}[5][french]
\item \rep{a}
\item \rep{b}
\item \rep{c}
\item \rep{b}
\item \rep{a}
\item \rep{c}
\item \rep{b}
\item \rep{c}
\item \rep{a}
\item \rep{a}
\end{reponses}
\pourquoi{question 6 : l'adjectif de couleur se place avant le nom et reste invariable.
Question 10 : \en{the umbrella} \phon{ði}, mais \en{the pen} \phon{ðə}.}
\end{corrigetest}
